\documentclass[11pt,oneside]{article}

\usepackage[letterpaper,margin=0.82in,headheight=15pt,footskip=30pt]{geometry}
\usepackage[T1]{fontenc}
\usepackage[utf8]{inputenc}
\usepackage{mathpazo}
\usepackage[scaled=0.92]{helvet}
\usepackage{courier}
\usepackage{microtype}
\usepackage{amsmath,amssymb}
\usepackage{xcolor}
\usepackage{parskip}
\usepackage{setspace}
\usepackage{enumitem}
\usepackage{titlesec}
\usepackage{fancyhdr}
\usepackage{lastpage}
\usepackage{hyperref}
\usepackage{xurl}
\urlstyle{same}
\usepackage{bookmark}
\usepackage[most]{tcolorbox}
\usepackage{ragged2e}
\usepackage{etoolbox}
\usepackage{needspace}
\usepackage{endnotes}

\definecolor{ManifestInk}{HTML}{202124}
\definecolor{ManifestGray}{HTML}{5F6368}
\definecolor{ManifestLight}{HTML}{F3F4F5}
\definecolor{ManifestRule}{HTML}{8A8D91}

\hypersetup{
  colorlinks=true,
  linkcolor=ManifestInk,
  urlcolor=ManifestGray,
  pdftitle={The Ethical Quiet Quitting Manifesto},
  pdfauthor={Leah},
  pdfsubject={Repair First. Stop Hiding the Damage. Withdraw Only When
Repair Fails.},
  pdfkeywords={quiet quitting, ethics, work, labor, management, safety, employee voice}
}

\setstretch{1.04}
\setlength{\parindent}{0pt}
\setlength{\parskip}{0.50em}
\setcounter{secnumdepth}{-1}
\setcounter{tocdepth}{2}
\emergencystretch=3em
\clubpenalty=10000
\widowpenalty=10000
\displaywidowpenalty=10000

\titleformat{\section}
  {\Large\sffamily\bfseries\color{ManifestInk}}
  {}{0pt}{}
  [\vspace{0.2em}\color{ManifestRule}\titlerule]
\titlespacing*{\section}{0pt}{2.2em}{0.8em}

\titleformat{\subsection}
  {\large\sffamily\bfseries\color{ManifestInk}}
  {}{0pt}{}
\titlespacing*{\subsection}{0pt}{1.5em}{0.45em}

\pretocmd{\section}{\Needspace{6\baselineskip}}{}{}
\pretocmd{\subsection}{\Needspace{4\baselineskip}}{}{}

\setlist[itemize]{leftmargin=1.45em,itemsep=0.22em,topsep=0.35em}
\setlist[enumerate]{leftmargin=1.65em,itemsep=0.32em,topsep=0.4em}
\providecommand{\tightlist}{%
  \setlength{\itemsep}{0.22em}\setlength{\parskip}{0pt}}

\renewenvironment{quote}
  {\begin{tcolorbox}[
      enhanced,
      breakable,
      colback=ManifestLight,
      colframe=ManifestRule,
      boxrule=0pt,
      leftrule=3pt,
      arc=0pt,
      outer arc=0pt,
      left=10pt,right=10pt,top=8pt,bottom=8pt,
      before skip=0.8em,after skip=0.8em
    ]\itshape}
  {\end{tcolorbox}}

\pagestyle{fancy}
\fancyhf{}
\fancyhead[L]{\sffamily\footnotesize\color{ManifestGray}The Ethical
Quiet Quitting Manifesto}
\fancyhead[R]{\sffamily\footnotesize\color{ManifestGray}Leah}
\fancyfoot[C]{\sffamily\footnotesize\color{ManifestGray}\thepage\ / \pageref*{LastPage}}
\renewcommand{\headrulewidth}{0.3pt}
\renewcommand{\footrulewidth}{0pt}

\let\footnote=\endnote
\renewcommand{\notesname}{Notes}

\begin{document}
\hypersetup{pageanchor=false}

\begin{titlepage}
\thispagestyle{empty}
\vspace*{0.9in}

{\sffamily\bfseries\fontsize{31}{36}\selectfont\color{ManifestInk}
The Ethical Quiet Quitting Manifesto\par}

\vspace{0.5em}
{\color{ManifestRule}\rule{\textwidth}{1.2pt}}
\vspace{1.0em}

{\sffamily\fontsize{17}{23}\selectfont\color{ManifestGray}
Repair First. Stop Hiding the Damage. Withdraw Only When Repair
Fails.\par}

\vfill

\begin{tcolorbox}[
  enhanced,
  colback=ManifestLight,
  colframe=ManifestRule,
  boxrule=0pt,
  leftrule=4pt,
  arc=0pt,
  left=14pt,right=14pt,top=13pt,bottom=13pt
]
{\Large\itshape We withdraw the subsidy, not the integrity.}
\end{tcolorbox}
\vfill

{\sffamily\large\bfseries Leah\par}
\vspace{0.2em}
{\sffamily\color{ManifestGray}August 16, 2026 \quad Version 1.5.4\par}
\vspace{0.55in}
\end{titlepage}

\pagenumbering{roman}
\thispagestyle{plain}
\tableofcontents
\clearpage
\pagenumbering{arabic}
\hypersetup{pageanchor=true}

\hypertarget{preamble}{%
\section{Preamble}\label{preamble}}

\textbf{Quiet quitting proper is not outright refusal of labor.} It
retains the defensible baseline of the job and withdraws only the
unowned excess demanded above it. Whether and when outright refusal is
ethical is a separate question this manifesto does not purport to
settle.

Work is a cooperative system, but employment is not a meeting of equal
power. The worker's most immediate remaining power is limited but real:
how they perform the paid role and whether they continue donating
excess. Economic necessity constrains even that choice. The employer
controls pay, staffing, schedules, rules, discipline, and much of the
field in which that labor must act.

That asymmetry creates different ethical ideals and different burdens.
Here, an \textbf{ideal} is an orienting ethical standard, not an
optional courtesy, a description of every real workplace, or a claim
that the burdens are equal. The worker-side ideal is the conscious and
responsible use of the limited power the worker actually exercises:
competent, safe, honest labor within the real job, and ownership of
chosen acts within what the worker can reasonably know and control. The
employer-side ideal is humane restraint and affirmative repair
proportionate to its greater control: sufficient resources, intelligible
priorities, usable feedback channels, and some plausible connection
between effort and a self-directed future. These are not mirror images.
Greater authority, information, and reach create the larger affirmative
burden and the wider consequences. For those who hold concentrated
power, humane restraint is not pathological tolerance of deception,
serious harm, or deliberate destruction. It is the refusal to use
concentrated power carelessly or punitively. No worker owes unlimited
self-erasure to keep a broken arrangement looking functional.

No system is perfect, and neither ideal requires perfection. Ordinary
friction, a hard week, a minor defect, or an isolated repairable mistake
does not by itself create the conditions for ethical quiet quitting
proper. At minimum, quiet quitting proper requires a material and
persistent mismatch between the defensible job and operative
expectations, and pathways toward meaningful correction must be absent,
unsafe, reasonably judged futile, unreasonably costly, or limited to
assimilation rather than repair; the further ethical conditions below
still apply. Imperfection calls for proportionate repair. Structural
closure is what makes withdrawal the last boundary.

Quiet quitting proper should not be the first move. It should be the
last honest boundary after good-faith repair has failed.

\begin{quote}
\textbf{The central rule:} Try to make the workplace better in every
ethical way reasonably available. When those pathways have been tested
and closed, stop using private sacrifice to counterfeit institutional
health. Perform the real job safely, honestly, and sustainably. That is
quiet quitting proper.
\end{quote}

Repair first is not permission-seeking. It helps the worker learn what
is true before exercising power and preserves the causal story of what
follows. It does not give the employer a right to endless chances, force
the worker through danger or retaliation, or transfer responsibility for
institutional repair to the person with the least control over it.

This argument finally crystallized through one workplace: chronic
under-capacity, attempted repair, and the hidden effort that kept
incompatible demands looking compatible. I offer the pattern for
recognition and criticism, not as a template other workers must prove
they fit before they are allowed a boundary.

\hypertarget{what-quiet-quitting-is-and-what-it-is-not}{%
\section{1. What Quiet Quitting Is, and What It Is
Not}\label{what-quiet-quitting-is-and-what-it-is-not}}

This manifesto uses \textbf{quiet quitting proper} in a precise sense:

\begin{quote}
Quiet quitting proper means competently performing the work one is paid
and reasonably expected to perform, while ceasing the discretionary,
uncompensated, unsustainable, or identity-consuming excess that has
silently become part of the institution's operating model.
\end{quote}

In ordinary speech, however, \textbf{quiet quitting} is a much looser
label. It may be used for a worker reserving capacity for moonlighting
or a second job; setting an ordinary boundary; disengaging for reasons
unrelated to the structure of the job; or deliberately trying to do as
little as possible regardless of whether the bargain is fair. Those are
not one moral object. Some may be harmless or justified; some may
reflect a worker's tacit adaptation to the lived bargain; some may cross
into neglect, retaliation, or misconduct. A shared label does not make
the acts ethically identical, and the label alone settles none of them.

My question is narrower and begins with my own case: what does quiet
quitting proper look like when I am actively trying to be a good person
while doing it---to do good, avoid unnecessary harm, seek capital-T
Truth to the best of my understanding, and give everyone affected,
including the employer, a fair shake---while withdrawing the unowned
excess? A fair shake is not a fiction of equal footing; it means judging
each actor's choices by what they could reasonably know and control.
This is not a claim of moral superiority or a verdict on every worker
described by the broader label. It is an attempt to remain answerable to
reality and other people while refusing self-erasure: to withdraw the
subsidy without withdrawing integrity.

\textbf{A disclosure about how I arrived here.} I did not begin by
practicing this framework. At earlier jobs, I engaged in what the
broader public label calls quiet quitting in ways I would not now call
quiet quitting proper. At times I treated my own needs as settling the
whole question. The mistake was not refusing self-sacrifice. It was
failing to test what I thought was true against reality, trace
foreseeable effects on other people, or distinguish withdrawing an
unowned subsidy from simply putting myself first. I do not offer this
manifesto to retroactively certify those choices; some of them do not
survive its test.

I kept returning to withdrawal despite real trepidation because
\emph{something felt right about quiet quitting that I couldn't put my
finger on}. That feeling was a clue, not a verdict. Reflection on the
consequences---and on the distance between what I thought was true and
what I could later justify as true---showed me both parts of the
problem: a worker can be responding to a real structural mismatch and
still respond to it unethically.

This manifesto is the result of forcing that unresolved intuition
through truth-testing, consequence-awareness, and responsibility. My
responsibility for the acts I chose does not put the parties on equal
footing or erase the institution's larger structural burden. This
disclosure explains my route to the framework. It is not a confession
other workers owe before the facts of their boundary can be recognized.

Except where the broader public label is explicitly under discussion,
\textbf{quiet quitting} below is shorthand for \textbf{quiet quitting
proper}.

For the ethical quiet quitter proper, the relation runs one way:

\[
\boxed{\text{quiet quitter proper}\ \Longrightarrow\ \text{structural or managerial expectations exceed the defensible job baseline}.}
\]

The baseline is not the least an employer can compel without
consequence. It is the compensated, reasonably defined role---including
legitimate professional and safety standards---performed competently at
a sustainable pace. Quiet quitting proper begins when structural or
managerial expectations exceed that baseline and the worker stops
supplying the unowned excess.

That excess condition is necessary, not sufficient. The pathway and
ethical conditions below distinguish quiet quitting proper from ordinary
boundary-setting, misconduct, or collective pressure.

It is not refusing to work. It is refusing to donate an undefined second
job.

No worker must earn permission to stop making a genuinely voluntary
gift. An institution does not acquire a new moral right to notice,
staging, or handoff merely because it organized itself around labor it
never sustainably resourced. Where a boundary may touch other people,
responsible power means seeing what can be seen and owning what one
actually chooses; it does not mean accepting responsibility for
conditions controlled elsewhere.

Quiet may be spoken or unannounced. Where openness is safe and useful, a
worker can state the boundary plainly. Where openness invites
retaliation, silence is not automatically dishonesty. What ends is the
donation of unowned excess, not the worker's integrity.

The real job is the work paid for and reasonably expected under the role
and its professional standards, not every task normalized through hidden
sacrifice. Repetition does not turn an unsustainable subsidy into a
permanent duty.

Simply declining a genuinely optional extra is ordinary
boundary-setting. It can be wholly ethical without being the narrower
phenomenon named here. Quiet quitting proper becomes distinct when
operative structural or managerial expectations exceed the defensible
job baseline. The gap may appear in the words---the codified role is so
vague that no defensible boundary can be identified, or is internally
contradictory; in the practice---the real job materially extends or
contradicts the written one without agreement, resources, or reciprocal
adjustment; or in the operating conditions---volume, pace, staffing,
equipment, or priorities make competent, sustainable performance of the
nominal bundle impossible. An ``optional'' extra enforced through
pressure, evaluation, or discipline is not optional in practice.

Correction clarifies the role, reconciles and prioritizes conflicting
standards, resources the work, or negotiates expanded duties the worker
can actually reject. Assimilation says: absorb the contradiction, supply
the missing capacity, and let hidden compensating labor become the job.

\begin{quote}
\textbf{Assimilation is not correction. It is the defect becoming the
worker's identity and duty.}
\end{quote}

In a healthy, properly resourced workplace, a worker performing this
same baseline conduct simply looks average: the job is being done
competently at a sustainable pace. There is nothing there to call quiet
quitting. In a broken workplace, the same conduct looks like withdrawal
only because the institution had mistaken excess for capacity.

\begin{quote}
\textbf{The behavior is ordinary work. The phenomenon is the institution
discovering that it had mistaken a subsidy for the job.}
\end{quote}

It is not:

\begin{itemize}
\tightlist
\item
  sabotaging output;
\item
  deliberately working below a reasonable professional standard;
\item
  falsifying records, hiding urgent risks, or manufacturing failures;
\item
  concealing serious safety risks or manufacturing harm and calling it
  ordinary quiet quitting;
\item
  secretly creating bottlenecks that would not otherwise exist and
  presenting them as ordinary reduced effort;
\item
  retaliating against a decent direct supervisor for failures above
  their authority;
\item
  performing helplessness as a strategy;
\item
  proof that the worker rejects ambition, discipline, or difficult work;
\item
  confusing revenge with a boundary.
\end{itemize}

It is:

\begin{itemize}
\tightlist
\item
  working at a sustainable professional pace;
\item
  fulfilling the defined role without silently absorbing every missing
  role around it;
\item
  declining chronic unpaid availability and unbounded role expansion;
\item
  refusing to turn every managerial failure into a personal emergency;
\item
  making incompatible demands visible instead of privately reconciling
  them through exhaustion;
\item
  preserving energy, time, health, and agency for a life beyond the
  institution;
\item
  reserving exceptional effort for work, people, and futures that can
  meaningfully receive it.
\end{itemize}

The ethical distinction is simple:

\[
\boxed{\text{Ceasing to subsidize a broken system} \neq \text{manufacturing damage to a system}.}
\]

That distinction defines \textbf{individual quiet quitting proper}. It
does not decide the morality of strikes, boycotts, slowdowns,
coordinated work-to-rule, or other collective action. Collective power
can intentionally create pressure and disruption; its ethics are
addressed in Section 9.

The warnings are different because the power is different. For workers,
the label cannot turn bad faith, retaliation, deliberate neglect, or
conduct that belongs to another category into ethical action; where the
act is consciously chosen, name it honestly and own it within what can
reasonably be known and controlled. For employers, the existence of
unlike cases cannot erase the structurally produced one or justify
treating every withdrawal as laziness or misconduct. Because employers
control classification, discipline, staffing, and repair, they bear the
larger duty to distinguish the cases fairly before acting.

\begin{quote}
\textbf{A workplace can teach quiet quitting before a worker ever calls
it that.}
\end{quote}

This makes no statistical claim about what most people described by the
broader label do. My case is one explicitly articulated version of a
structural possibility developed in the next section, not proof that
less articulated cases lack reasons or integrity. If even this
consciously examined path reaches withdrawal, an institution cannot make
articulate self-explanation another condition of taking structurally
produced withdrawal seriously.

\hypertarget{why-workers-withdraw}{%
\section{2. Why Workers Withdraw}\label{why-workers-withdraw}}

In the narrower case developed here, withdrawal need not mean that the
worker stopped caring. Care may persist after actionable pathways to
repair have closed.

My working hypothesis is that many cases described by the broader public
label may share this structure without sharing this manifesto's explicit
reasoning. The posture can emerge gradually: the environment teaches a
person what effort leads to, while that person's circumstances and
values shape what they continue to give, what they protect, and what
they withdraw. They may settle into the mindset without ever formalizing
it as an ethical program. That is a proposed explanation, not a measured
prevalence claim.

This manifesto is not a defense of apathy. It addresses workers who try
to repair what they reasonably can, tell the truth as far as they can,
and eventually stop feeding a system that converts every rescue into
evidence that no rescue was needed.

The worker learns, repeatedly, that:

\[
\text{speaking} \nrightarrow \text{change},
\]

\[
\text{excellent work} \nrightarrow \text{meaningful advancement},
\]

\[
\text{identifying a problem} \nrightarrow \text{repair},
\]

\[
\text{sacrifice} \rightarrow \text{the new minimum expectation}.
\]

A promotion may exist in name but not in meaning. Requests for help may
be acknowledged but not acted upon. The chain of command may transmit
orders downward while filtering inconvenient reality upward. The
employee may still possess the physical ability to act, but the routes
connecting action to a self-directed future have narrowed or closed.

That is not merely low morale. It is a contraction of practical agency.

Yet perceived closure can be mistaken, temporary, or caused by an
ordinary communication failure.\footnote{Elizabeth W. Morrison,
  ``Employee Voice and Silence,'' \emph{Annual Review of Organizational
  Psychology and Organizational Behavior} 1 (2014): 173-197,
  \url{https://doi.org/10.1146/annurev-orgpsych-031413-091328}. The
  review describes organizational and individual conditions that can
  encourage or inhibit employee voice; it does not establish that every
  apparently closed pathway is in fact closed.} That is why ethical
quiet quitting begins by testing the closure where doing so is
reasonably safe and useful. The test serves the worker's own judgment.
It is not a presumption that the institution deserves another appeal.

\hypertarget{the-exit-fallacy-why-not-just-get-another-job}{%
\subsection{2.1 The Exit Fallacy: ``Why Not Just Get Another
Job?''}\label{the-exit-fallacy-why-not-just-get-another-job}}

The standard objection to quiet quitting is: ``If you can still do the
job properly, why not keep doing it here - or just get another job?''

It sounds like common sense only because it assumes exit is immediate,
costless, and safe.

\begin{quote}
\textbf{``Why not just get another job?'' is often the labor-market
version of ``If they cannot afford bread here, let them eat cake.''} It
answers a constraint by assuming access to the very resources the
constraint has already taken away.
\end{quote}

The comparison is deliberate hyperbole. It does not equate the suffering
involved, and it makes no historical claim about who coined the phrase.
It identifies the same error of answering a constraint by recommending
an option the constrained person may not actually be able to reach.

People need jobs to live. Wages keep housing, food, medicine, utilities,
transportation, and the rest of ordinary survival available. For workers
without enough financial cushion to absorb an income gap, missing even
one paycheck can produce immediate material harm. A job search itself
consumes scarce resources: time, attention, internet access,
transportation, scheduling flexibility, references, emotional bandwidth,
and tolerance for uncertainty. A bad job can consume exactly the time
and mental clarity required to escape it.\footnote{Board of Governors of
  the Federal Reserve System, \emph{Economic Well-Being of U.S.
  Households in 2025} (Washington, DC, May 2026),
  \url{https://www.federalreserve.gov/publications/files/2025-report-economic-well-being-us-households-202605.pdf};
  West Health and Gallup, ``One in Four U.S. Employees Locked in Jobs
  for Health Insurance,'' July 21, 2026,
  \url{https://news.gallup.com/poll/712166/one-four-employees-locked-jobs-health-insurance.aspx};
  Daniel Schneider and Kristen Harknett, ``Consequences of Routine
  Work-Schedule Instability for Worker Health and Well-Being,''
  \emph{American Sociological Review} 84, no. 1 (2019): 82-114,
  \url{https://doi.org/10.1177/0003122418823184}. These sources document
  forms of financial fragility, benefit dependence, and schedule
  instability in defined U.S. samples. They do not establish why any
  particular worker remains in a job or that the schedule study is
  nationally representative.}

In agency terms, another job may be formally possible without being a
live option:

\[
\mathcal E_t^{\mathrm{formal}}\neq\varnothing
\quad\nRightarrow\quad
\mathcal E_t^{\mathrm{viable}}\neq\varnothing.
\]

A vacancy may exist somewhere. That does not mean this worker can reach
it, be hired, survive the gap before the first paycheck, accept its
schedule, replace lost benefits, or risk losing current income while
applying. Thus:

\[
\text{another job exists}
\nRightarrow
\text{another job is presently reachable},
\]

and

\[
\text{ability to perform this job}
\nRightarrow
\text{ability to survive an employment transition}.
\]

The current employer can also consume the worker's exit capacity through
overtime, unstable scheduling, exhaustion, and chronic crisis. Telling
that worker to leave can amount to demanding that they finance the
employer's failure with unemployment risk.

Quiet quitting proper can therefore be a bridge, not merely a
destination. By withdrawing unsustainable excess while preserving
competent core work and income, a worker may recover enough time,
health, and cognitive capacity to search, train, organize, or build a
genuinely live exit. In this role, quiet quitting is configurational
action: it can create the conditions under which leaving becomes
possible. It can also bring discipline, retaliation, lost advancement,
or dismissal. Reclaimed power is real power, with real risk.

None of this means a worker should never leave. A genuinely live better
option should be pursued when appropriate. The point is that its
liveness must be assessed, not presumed. Continued attendance does not
prove consent, satisfaction, or endorsement. It may mean only that rent
is due.

The wage relation does not produce a clean binary between fully invested
and gone. One important intermediate state is a person who still needs
the wage but can no longer rationally donate hidden heroics. That is
precisely the terrain of quiet quitting. If management answers that
state with more pressure rather than restored pathways, ordinary mistake
risk and forced triage can intensify together.

\hypertarget{grind-culture-hard-work-needs-somewhere-to-go}{%
\subsection{2.2 Grind Culture: Hard Work Needs Somewhere to
Go}\label{grind-culture-hard-work-needs-somewhere-to-go}}

This manifesto is not an attack on ambition, discipline, or hard work.
Deep effort can be a rational joy when it builds a future the worker can
recognize as their own. There are jobs, crafts, teams, missions, and
businesses for which long hours and difficult work can be rational,
joyful, and personally meaningful. A person may willingly grind because
the effort produces mastery, security, ownership, belonging, influence,
credible advancement, or a future they genuinely recognize as their own.

That is good. Someone who has found work that resonates with them should
be glad they found it.

The error begins when a contingent fit is converted into a universal
character test:

\begin{quote}
\textbf{I can grind here, therefore anyone who does not grind here is
lazy, weak, entitled, or morally unserious.}
\end{quote}

Meaningful grind is not generated by attitude alone. It is co-produced
by a worker, a role, a manager, a team, a reward structure, a life
stage, and a particular moment in the world. Access to it can depend
heavily on chance: meeting the right people, entering the right field at
the right time, having enough stability to take a risk, finding
coworkers whose own versions of commitment are compatible, and landing
somewhere that can actually receive what one is capable of giving.

A person can search for that alignment, cultivate skills, take risks,
and improve the odds. They cannot unilaterally command a workplace to
become meaningful, force a manager to become trustworthy, create
advancement where none exists, or make everyone around them participate
in the same reciprocal project. The right attitude cannot, by itself,
manufacture the right social and material field.

Schematically, discretionary grind depends on more than character:

\[
G_t^{\mathrm{disc}}
=
f(R_t,O_t,Q_t,P_t,C_t,B_t),
\]

where:

\begin{itemize}
\tightlist
\item
  \(R_t\) is personal resonance with the work;
\item
  \(O_t\) is ownership, autonomy, or meaningful influence;
\item
  \(Q_t\) is tangible reciprocal return;
\item
  \(P_t\) is a credible pathway from present effort to a self-directed
  future;
\item
  \(C_t\) is the surrounding field of compatible coworkers, leadership,
  and institutional support;
\item
  \(B_t\) is the worker's available health, time, and cognitive
  bandwidth.
\end{itemize}

This is not a psychometric formula. It is a reminder that exceptional
effort is an ecological outcome, not a scalar reading of moral worth.

Grind ideology often collapses the whole vector into one flattering
verdict:

\[
\text{high discretionary effort}
\approx
\text{strength, seriousness, and virtue}.
\]

But the valid implications are much weaker:

\[
\begin{aligned}
\text{high discretionary effort}
&\nRightarrow
\text{superior character},\\
\text{withdrawal of excess}
&\nRightarrow
\text{weakness or absence of ambition}.
\end{aligned}
\]

The same person may grind relentlessly for their own company, their art,
their family, a team they trust, a craft they love, or an employer that
tangibly shares the return---and quiet quit at an institution where
extraordinary effort merely disappears into the next quarter's minimum
expectation. That is not necessarily inconsistency. The objects are
different.

\begin{quote}
\textbf{I would grind for a company if the grind tangibly built
something I could recognize as mine: security, mastery, influence,
belonging, ownership, a credible future, or a mission I actually share.
I will not grind merely to prove that I am hard, that I work hard, that
I am not weak, or that I can endure being ground down.}
\end{quote}

When tangible return is absent, grind culture can substitute an
\textbf{identity wage}: ``I am hard. I work hard. I am not weak. I am a
grinder.'' Alison R. Buck uses \emph{identity wages} in her study of
game-industry labor.\footnote{Alison R. Buck, ``The Price of Consent:
  Identity Wages in the Games Industry,'' \emph{Journal of Contemporary
  Ethnography} 51, no. 6 (2022): 868-894,
  \url{https://doi.org/10.1177/08912416221085558}.} I adapt the term
here for a narrower warning beyond that domain. That identity may be
personally chosen and genuinely valuable to someone. But it cannot be
treated as compensation owed by the worker to the employer. If suffering
itself becomes proof of virtue, then exploitation becomes self-sealing:
every injury demonstrates toughness, every boundary demonstrates
weakness, and no evidence can show that the arrangement is a bad
bargain.

Hard work does not need to be pleasant every hour to be worthwhile.
Meaningful projects include boredom, frustration, repetition, sacrifice,
and failure. Nor does personal non-resonance erase the worker's core
obligations after accepting wages. The distinction concerns
\textbf{discretionary excess}: the unbounded energy above competent,
safe, honest performance.

Quiet quitting proper is therefore not necessarily anti-work or
anti-ambition. It can be the refusal to let one institution monopolize a
person's highest effort without providing a credible reciprocal basis.
The recovered capacity may go into education, art, caregiving,
organizing, recovery, a job search, a future business, or another
workplace capable of receiving it.

\begin{quote}
\textbf{Hard work becomes investment when it builds an owned future.
Without that connection, the same effort becomes depletion.}
\end{quote}

This does not make every difficult demand illegitimate. Workers should
not dismiss one merely because it does not feel inspiring. But managers
must not use the most personally resonant employee as the neutral
baseline for everyone else. Passion is real, powerful, and partly
contingent. It is not payroll infrastructure.

The question of whether hard work still builds an owned future has a
specifically American form.

\hypertarget{the-hard-working-american-and-the-american-dream}{%
\subsection{2.3 The Hard-Working American and the American
Dream}\label{the-hard-working-american-and-the-american-dream}}

In the strand of American culture examined here, the
\textbf{hard-working American} and the \textbf{American Dream} are not
independent ideas. They are the input and output sides of one social
promise.\footnote{Library of Congress, ``The American Dream,'' classroom
  materials, accessed August 16, 2026,
  \url{https://www.loc.gov/classroom-materials/american-dream/}; Raj
  Chetty et al., ``The Fading American Dream: Trends in Absolute Income
  Mobility Since 1940,'' \emph{Science} 356, no. 6336 (2017): 398-406,
  \url{https://doi.org/10.1126/science.aal4617}; Gabriel Borelli,
  ``Americans Are Split over the State of the American Dream,'' Pew
  Research Center, July 2, 2024,
  \url{https://www.pewresearch.org/short-reads/2024/07/02/americans-are-split-over-the-state-of-the-american-dream/}.
  Chetty and colleagues estimate that roughly 90 percent of the 1940
  birth cohort had higher inflation-adjusted household incomes than
  their parents at about age thirty, compared with roughly half of
  people born in the 1980s. The materials do not reduce the Dream to
  income, establish equal access, or prove Leah's mechanism. Her
  input/output account is one influential strand; ``all input and no
  output'' refers to exceptional sacrifice \textbf{here}, not the
  literal absence of wages or present benefits.}

\begin{quote}
\textbf{The hard-working American supplies the effort. The American
Dream is the future that effort is supposed to make reachable.}
\end{quote}

The input-side ideal praises discipline, reliability, endurance,
initiative, contribution, and the willingness to sacrifice now for
something larger later. The output-side ideal says that such effort can
plausibly become a stable life: enough material security to breathe, a
home or other durable stake in the world, the ability to support a
family, education and opportunity, retirement, ownership, upward
movement, independence, dignity, and a future increasingly shaped by the
worker rather than merely endured by them.

The promise has never meant that every hardworking person is guaranteed
every desired outcome. Chance, starting position, health, geography,
timing, social connection, discrimination, family obligations, and the
actions of other people all affect what becomes reachable. But the ideal
retains moral force only while the causal arrow remains credible:

\[
\text{hard work}
\longrightarrow
\text{a materially and personally owned future}.
\]

A worker can tolerate uncertainty in that arrow. They cannot rationally
invest forever after it has become indistinguishable from fiction.

When wages cannot support a stable life, schedules destroy the worker's
ability to plan one, promotions change titles without changing futures,
benefits remain precarious, ownership never increases, and every burst
of productivity is converted into the next permanent minimum, the
worker's own estimate may begin to change. Let \(\widehat P_t\) denote
that judgment; schematically, not as a fitted population estimate:

\[
\widehat P_t(\text{owned future}\mid\text{additional effort here})\downarrow.
\]

The critical word is \textbf{here}. The worker may still believe in hard
work. They may still want the house, the family, the craft, the
business, the security, the independence, or the chance to give their
children more. What has collapsed is the belief that additional
sacrifice for this institution is a live route toward those goods.

At that point, appeals to the ``hard-working American'' can become an
identity wage dressed in national mythology. The institution continues
demanding the input after severing it from the promised output:

\begin{quote}
Work harder because hardworking people are good people. Endure because
endurance proves character. Be proud that you can carry what we refuse
to repair.
\end{quote}

But a social bargain cannot remain coherent as a one-sided character
demand. An employer cannot substitute borrowed national ideals for pay,
security, ownership, advancement, or meaningful influence. Patriotism is
not compensation, and suffering is not evidence that a future is being
built.

\begin{quote}
\textbf{The American Dream cannot function as deferred compensation
after the pathway from present labor to that future has ceased to be
credible.}
\end{quote}

This is one reason quiet quitting follows pathway closure. The worker
does not necessarily reject the Dream. They cease treating the employer
as its vehicle. They preserve competent core work because survival still
depends on the wage, while withdrawing the discretionary sacrifice that
no longer appears to purchase any owned future.

Schematically:

\[
\text{dream-path closure}
\longrightarrow
\text{withdrawal of discretionary investment}.
\]

Quiet quitting can therefore preserve the work ethic rather than destroy
it. The reclaimed effort may be redirected toward education, recovery,
family, art, organizing, a second income, a job search, a future
business, or another institution where sacrifice and return remain
causally connected. The person is not necessarily refusing to build.
They are refusing to keep building a future that belongs entirely to
someone else.

The proportionality rule remains important. The claim is not that one
ordinary job guarantees wealth, fulfillment, ownership, or limitless
upward mobility. It is that the more an institution invokes the
hard-working-American ideal to demand exceptional loyalty, sacrifice,
flexibility, and output, the stronger its obligation to make the return
tangible and the future pathway credible.

\begin{quote}
\textbf{The American work ethic without a reachable American Dream is
all input and no output. At that point it is no longer a reciprocal
ideal; it is extraction wearing the language of character.}
\end{quote}

\hypertarget{repair-first-the-ethical-order-of-operations}{%
\section{3. Repair First: The Ethical Order of
Operations}\label{repair-first-the-ethical-order-of-operations}}

Where repair remains safe and possible, quiet quitting is strongest when
it follows an honest attempt rather than an untested assumption. That is
truth-seeking, not a permission ritual.

\hypertarget{diagnose-before-accusing}{%
\subsection{3.1 Diagnose before
accusing}\label{diagnose-before-accusing}}

First determine what can reasonably be known about what is happening.

Distinguish a structurally impossible workload from a temporarily
difficult week. Separate what can be seen of understaffing, bad
scheduling, missing equipment, poor sequencing, unclear priorities,
training gaps, demand spikes, or simple misunderstanding. Notice where
the real constraint appears to live and who has authority to investigate
and change it.

Do not mistake suspicion for proof. Reconstruct what you reasonably can;
the institution controls the fuller record.

The question is not merely, ``Why am I suffering?'' It is:

\[
\text{What combination of demand, staffing, equipment, policy, and priority produces this outcome?}
\]

\hypertarget{use-the-nearest-safe-and-functioning-channel}{%
\subsection{3.2 Use the nearest safe and functioning
channel}\label{use-the-nearest-safe-and-functioning-channel}}

Start with the people who can still hear you without making the danger
worse.

Where that route is safe and useful, communicate with the direct
supervisor. Explain the problem in concrete terms. Separate facts from
interpretations. Describe the current resources, the demanded outcomes,
the conflict between them, and the smallest realistic repair.

A good direct supervisor is not merely another rung in a hierarchy. They
can become a local witness to reality, a translator across institutional
cultures, and a source of provenance when higher levels would otherwise
see only an unexplained bad number.

That may not be the available route. A worker need not exhaust a channel
reasonably expected to retaliate, discriminate, expose protected
information, or create serious danger. Repair first asks for a
good-faith attempt where a safe and functioning attempt exists---not
ritual vulnerability before a hierarchy that has already made
truth-telling hazardous.

\hypertarget{propose-repair-not-merely-distress}{%
\subsection{3.3 Propose repair, not merely
distress}\label{propose-repair-not-merely-distress}}

Distress matters, but institutions often know how to ignore distress.
They have more difficulty ignoring a clearly specified constraint paired
with a feasible intervention.

Where possible, propose actions such as:

\begin{itemize}
\tightlist
\item
  adding labor at a specific bottleneck or time window;
\item
  changing a production sequence;
\item
  reducing or reordering lower-priority work;
\item
  repairing or replacing a capacity-limiting tool;
\item
  clarifying which standard takes precedence when standards conflict;
\item
  creating an escalation rule for predictable overload;
\item
  changing schedules to match observed demand rather than idealized
  demand.
\end{itemize}

The goal is not to do management's entire job for free. Offering one
realistic intervention can test whether a functioning pathway exists
without making the worker responsible for building it.

\hypertarget{force-incompatible-priorities-into-the-open}{%
\subsection{3.4 Force incompatible priorities into the
open}\label{force-incompatible-priorities-into-the-open}}

When the available resources cannot satisfy all demands, do not
privately absorb the contradiction.

Say, in substance:

\begin{quote}
With the present staffing and demand, we can complete A and B safely, or
A and C safely, but not A, B, and C at the demanded standard. Which
takes priority?
\end{quote}

This is one of the cleanest ethical instruments available to a worker.
It transfers the tradeoff to the level that owns the authority to make
it.

When management chooses a lawful and safe priority, follow it honestly.
The unchosen work is then not an unexplained personal failure. It is the
visible consequence of an explicit allocation decision. A priority order
cannot erase the floor of safety, legality, honesty, or core
professional integrity.

\hypertarget{escalate-proportionally}{%
\subsection{3.5 Escalate proportionally}\label{escalate-proportionally}}

Use the chain of command while it remains a communication channel. When
it becomes only a shield against unwelcome information, use other
ethical and available routes: documented escalation, established
complaint systems, collective worker discussion, union representation
where applicable, or external reporting for genuine safety or legal
concerns.

Escalation should be truthful, proportional, and specific. It should aim
at correction rather than humiliation.

\hypertarget{communicate-in-a-channel-the-institution-can-read}{%
\section{4. Communicate in a Channel the Institution Can
Read}\label{communicate-in-a-channel-the-institution-can-read}}

Workers often discover that explanations are treated as opinions while
money, labor hours, waste, overtime, throughput, and customer complaints
are treated as reality.

That does not mean money is the message. Money is only a channel.

A negative number by itself is radically ambiguous. It may be decoded as
understaffing, incompetence, poor management, bad equipment, unusual
demand, laziness, or random variation. A financial consequence without
causal provenance can communicate the opposite of what the worker
intended.

The useful signal is not:

\[
\Delta \$ < 0.
\]

It is:

\[
(\text{staffing},\ \text{demand},\ \text{equipment},\ \text{priority decision})
\longrightarrow
\text{observable outcome}.
\]

The arrow records a causal hypothesis for the institution to
investigate, not proof that these variables exhaust the cause.

A consequence communicates most clearly when whatever source-map is
safely available remains attached.

Where it is safe, lawful, and useful, preserve an honest record of
staffing levels, workload, material shortages, equipment failures,
explicit priority decisions, and the resulting backlog or cost. Respect
privacy and do not remove confidential or protected information merely
to build a case. Do not falsify, exaggerate, or stage anything. The
worker need not prove a unique cause. Management, with broader access to
system data and control over resources, owns the deeper causal
investigation.

The principle is:

\begin{quote}
Do not create a failure in order to send a message. Stop erasing the
evidence of a failure that already exists.
\end{quote}

\hypertarget{ethical-de-buffering-stop-counterfeiting-capacity}{%
\section{5. Ethical De-Buffering: Stop Counterfeiting
Capacity}\label{ethical-de-buffering-stop-counterfeiting-capacity}}

A workplace can appear to be functioning at the capacity management
believes it possesses while actually relying on an invisible worker
subsidy.

Schematically, let

\[
Y = F(D,N,E,H),
\]

where:

\begin{itemize}
\tightlist
\item
  \(Y\) is observed output;
\item
  \(D\) is demand;
\item
  \(N\) is staffing;
\item
  \(E\) is equipment and process capacity;
\item
  \(H\) is hidden heroic compensation: unsustainable pace, role
  absorption, skipped recovery, exceptional flexibility, constant
  rescue, and unrecorded cognitive or emotional labor.
\end{itemize}

When management treats output sustained by a large \(H\) as proof that
\(N\) and \(E\) were sufficient, it learns the wrong lesson. The
worker's competence becomes evidence against the worker's own request
for help.

\textbf{Ethical de-buffering} means returning \(H\) toward a sustainable
baseline while holding safety, honesty, and professional integrity
fixed.

It may mean:

\begin{itemize}
\tightlist
\item
  working quickly but not frantically;
\item
  taking agreed or required breaks;
\item
  declining to cover chronic vacancies as an indefinite personal
  obligation;
\item
  ending habitual off-clock problem-solving;
\item
  refusing to treat every predictable rush as an unforeseeable
  emergency;
\item
  allowing noncritical backlog to remain visible;
\item
  reporting unfinished lower-priority work instead of hiding it through
  unpaid effort;
\item
  requiring explicit authorization before taking on materially expanded
  duties;
\item
  preserving enough capacity to perform tomorrow's work rather than
  consuming tomorrow to save today.
\end{itemize}

Heroic effort may be freely chosen during a genuine emergency. It cannot
ethically be converted into the permanent baseline and then used to
prove that the emergency staffing level is adequate.

\begin{quote}
\textbf{The worker's body is not the slack variable in a staffing
model.}
\end{quote}

Ethical de-buffering is not making the system fail. It is refusing to
counterfeit success.

But hidden heroic compensation can be the last remaining margin between
an overloaded system and acute constraint collapse. If that margin is
withdrawn, the consequences may not scale smoothly. A small reduction in
hidden effort can move a brittle workplace across a threshold at which
delays, mistakes, corner-cutting, conflict, and backlog amplify one
another. That risk does not create an obligation to keep supplying the
subsidy forever. The worker remains responsible for safety and truth
within the choices they control. Management, which controls staffing,
process, and priorities, bears the primary duty to plan for the real
capacity that remains.

\hypertarget{the-fuck-it-factor-acute-pathway-collapse}{%
\section{6. The ``Fuck-It Factor'': Acute Pathway
Collapse}\label{the-fuck-it-factor-acute-pathway-collapse}}

A workplace can be structurally broken over months and still fail moment
by moment inside a single rush.

The model so far has described the long horizon: the worker speaks,
proposes, compensates, and waits, but the connection between action and
meaningful future outcome disappears. When the institution has also
folded unowned excess into the job, quiet quitting proper is the
long-horizon response: withdrawal of that excess after meaningful
correction pathways close.

That long-horizon account needs a moment-to-moment counterpart.

\begin{quote}
\textbf{The fuck-it factor}, as I use the term, is the felt transition
from trying to satisfy the whole constraint set to judging, under time
pressure and reduced mental clarity, that no available action can
satisfy it all while action is still required. Something must give, so
the person tries to sacrifice what appears to be the least important
thing.
\end{quote}

The deliberately blunt name preserves the experience as lived.

The fully compliant route may truly be absent, or it may only appear
absent within the time, information, equipment, and authority available.
The phenomenology is real in either case; which condition obtained is a
later factual question.

A rush at work can require one person to satisfy many constraints
simultaneously:

\begin{itemize}
\tightlist
\item
  preserve safety and sanitation;
\item
  preserve product or service quality;
\item
  meet queue and timing targets;
\item
  sequence scarce equipment correctly;
\item
  respond to customer-specific requirements;
\item
  prepare for the next wave of demand;
\item
  coordinate with coworkers;
\item
  obey managerial priorities;
\item
  regulate frustration and fear;
\item
  and act now rather than deliberate indefinitely.
\end{itemize}

Let the live constraints at time \(t\) be

\[
\mathcal C_t = \{c_1,c_2,\ldots,c_n\},
\]

and let \(\mathcal A_t\) be the actions presently available. The fully
compliant action-set is

\[
\mathcal A_t^{\mathrm{all}}
=
\{a\in\mathcal A_t : a\text{ satisfies every live constraint before its deadline}\}.
\]

Under sufficient demand, understaffing, equipment contention,
interruption, and time pressure, that set can become empty:

\[
\mathcal A_t^{\mathrm{all}}=\varnothing.
\]

The worker does not thereby gain the option to stop time. Orders
continue arriving. Food continues cooking. Customers continue waiting. A
manager still expects an answer. Some action must occur. Stopping,
slowing, refusing an unsafe instruction, escalating, or transferring
control are themselves actions. Safety is not merely another performance
target to trade against speed.

Descriptively---not as ethical permission to trade away safety---the
resulting decision is approximately:

\[
a_t^*
\approx
\operatorname*{arg\,min}_{a\in\mathcal A_t}
\sum_i \widehat w_{i,t}\,\ell_i(a),
\]

where \(\ell_i(a)\) is the loss created by violating constraint \(c_i\),
and \(\widehat w_{i,t}\) is the person's momentary estimate of how
important that constraint is. The hats matter. Under emotional strain,
divided attention, and severe urgency, those weights are being estimated
with reduced clarity.

Plainly:

\begin{quote}
\textbf{Fuck it. I cannot preserve everything. I will try to fuck the
least important thing.}
\end{quote}

This is not necessarily laziness, malice, or indifference. It names the
lived turn into triage, not a moral permission slip. The choice can
still be dangerous precisely because the person must choose while the
surrounding conditions may also make choosing well less likely.

\hypertarget{overload-departure-and-cascade}{%
\subsection{6.1 Overload, departure, and
cascade}\label{overload-departure-and-cascade}}

Overload, time pressure, stress, and interruption can raise ordinary
error risk in some tasks and conditions whether or not the fuck-it
factor appears.\footnote{James L. Szalma, Peter A. Hancock, and Sean
  Quinn, ``A Meta-Analysis of the Effect of Time Pressure on Human
  Performance,'' \emph{Proceedings of the Human Factors and Ergonomics
  Society Annual Meeting} 52, no. 19 (2008): 1513-1516,
  \url{https://doi.org/10.1177/154193120805201944}; Grant S. Shields,
  Matthew A. Sazma, and Andrew P. Yonelinas, ``The Effects of Acute
  Stress on Core Executive Functions: A Meta-Analysis and Comparison
  with Cortisol,'' \emph{Neuroscience \& Biobehavioral Reviews} 68
  (2016): 651-668,
  \url{https://doi.org/10.1016/j.neubiorev.2016.06.038}; Johanna I.
  Westbrook et al., ``Association of Interruptions with an Increased
  Risk and Severity of Medication Administration Errors,''
  \emph{Archives of Internal Medicine} 170, no. 8 (2010): 683-690,
  \url{https://doi.org/10.1001/archinternmed.2010.65}; Seyed M.
  Hashemian and Konstantinos Triantis, ``Production Pressure and Its
  Relationship to Safety: A Systematic Review and Future Directions,''
  \emph{Safety Science} 159 (2023): 106045,
  \url{https://doi.org/10.1016/j.ssci.2022.106045}. Together they
  support context-dependent risk, not a universal effect size or the
  coined factor. The factor is Leah's deliberately blunt
  phenomenological shorthand, not a validated scale, diagnosis, law,
  legal defense, or excuse.} The research establishes no universal
threshold or monotone law. The proposed second stage is Leah's
synthesis: conscious triage introduces another chance to choose the
wrong sacrifice or begin an unintended cascade. It may also avert a
worse outcome. Existing research does not isolate an independent
fuck-it-factor effect from the overload that produced it.

Three families must therefore remain distinct:\footnote{James Reason et
  al., ``Errors and Violations on the Roads: A Real Distinction?''
  \emph{Ergonomics} 33, nos. 10-11 (1990): 1315-1332,
  \url{https://doi.org/10.1080/00140139008925335}; Samuel J. Alper and
  Ben-Tzion Karsh, ``A Systematic Review of Safety Violations in
  Industry,'' \emph{Accident Analysis \& Prevention} 41, no. 4 (2009):
  739-754, \url{https://doi.org/10.1016/j.aap.2009.03.013}. The sources
  support distinguishing unintentional errors from intentional rule
  departures; neither implies that every intentional departure is
  malicious or mistaken.}

\begin{itemize}
\tightlist
\item
  \textbf{Intentional departures:} a person knowingly skips, delays,
  relaxes, or violates one requirement to preserve another. A departure
  is not automatically a mistake.
\item
  \textbf{Unintentional errors:} a person forgets a dependency, misreads
  the situation, executes poorly, or sacrifices the wrong constraint.
\item
  \textbf{Mixed cascades:} an intentional shortcut has an accidental
  consequence the person did not foresee or want.
\end{itemize}

Here, ``intentional'' does not mean that harm was desired. It means the
person knew some requirement was being abandoned because no fully
compliant route appeared live. A cascade can contain a deliberate
departure, an accidental error, or both at different points. Calling all
of them mistakes would erase the very distinction the model needs.

\hypertarget{repetition-can-normalize-the-workaround}{%
\subsection{6.2 Repetition can normalize the
workaround}\label{repetition-can-normalize-the-workaround}}

One fuck-it event can be emergency triage. Repeated, unexamined
workarounds can become familiar and harden into an operating regime when
the cause is never repaired.\footnote{Debra S. Debono et al., ``Nurses'
  Workarounds in Acute Healthcare Settings: A Scoping Review,''
  \emph{BMC Health Services Research} 13 (2013): 175,
  \url{https://doi.org/10.1186/1472-6963-13-175}; Anita L. Tucker, Amy
  C. Edmondson, and Steven Spear, ``When Problem Solving Prevents
  Organizational Learning,'' \emph{Journal of Organizational Change
  Management} 15, no. 2 (2002): 122-137,
  \url{https://doi.org/10.1108/09534810210423008}; National Aeronautics
  and Space Administration, \emph{NASA Human Factors Handbook},
  NASA-HDBK-8709.25, version 1.4 (2023),
  \url{https://standards.nasa.gov/sites/default/files/standards/NASA/Baseline/4/NASA-HDBK-870925-Baseline-14.pdf};
  Columbia Accident Investigation Board, \emph{Report, Volume I} (2003),
  \url{https://ntrs.nasa.gov/citations/20030093634}. These sources
  support conditional organizational processes of workaround persistence
  and normalization, not a mechanical decline in one person's threshold.}
Possible mechanisms include:

\begin{itemize}
\tightlist
\item
  the workaround becomes cognitively available and familiar;
\item
  the absence of an immediate disaster makes the breach feel safer than
  it was;
\item
  unfinished work and cleanup increase the next period's load;
\item
  exhaustion, resentment, and shame may reduce future adherence;
\item
  coworkers begin planning around the workaround;
\item
  management sees that the shift ``survived'' and treats the overloaded
  condition as proven capacity.
\end{itemize}

Feedback, rest, repair, training, accountability, or a bad outcome can
interrupt or reverse that pattern.

\hypertarget{the-short-horizon-dual-of-quiet-quitting}{%
\subsection{6.3 The short-horizon dual of quiet
quitting}\label{the-short-horizon-dual-of-quiet-quitting}}

The two processes have the same underlying shape at different
timescales:

\[
\text{long horizon: communicable and actionable pathways close}
\longrightarrow
\text{investment is withdrawn},
\]

\[
\text{short horizon: compliant operational pathways close under deadline}
\longrightarrow
\text{a constraint is sacrificed}.
\]

Quiet quitting proper is the long-horizon withdrawal of
institutionalized excess when both the future pathway and a meaningful
role-correction pathway have closed. The fuck-it factor is the acute,
moment-to-moment response to an action-space in which no fully compliant
move remains reachable before something must be done.

They can also reinforce each other. Chronic futility and resentment may
make acute abandonment more available. Acute mistakes, rule departures,
blame, and cleanup can further damage trust and close long-term
pathways. A workplace can therefore enter a feedback loop:

\[
\begin{aligned}
\text{under-capacity}
&\rightarrow
\text{acute triage}
\rightarrow
\text{mistakes and conflict}
\\
&\rightarrow
\text{less trust and less slack}
\rightarrow
\text{still greater under-capacity}.
\end{aligned}
\]

This loop is proposed, not inevitable.

\hypertarget{why-quiet-quitting-can-have-nonlinear-consequences}{%
\subsection{6.4 Why quiet quitting can have nonlinear
consequences}\label{why-quiet-quitting-can-have-nonlinear-consequences}}

Schematically, let effective momentary capacity be

\[
K_t = K_t^{\mathrm{base}} + H_t,
\]

where \(H_t\) is the hidden heroic compensation described above, and let
\(W_t\) be the load. The remaining slack is

\[
S_t = K_t - W_t.
\]

When \(H_t\) is withdrawn, \(S_t\) may cross from barely positive to
negative. Near saturation, some systems deteriorate nonlinearly rather
than in neat proportion.\footnote{J. F. C. Kingman, ``The Single Server
  Queue in Heavy Traffic,'' \emph{Proceedings of the Cambridge
  Philosophical Society} 57, no. 4 (1961): 902-904,
  \url{https://doi.org/10.1017/S0305004100036094}; David D. Woods, ``The
  Theory of Graceful Extensibility: Basic Rules That Govern Adaptive
  Systems,'' \emph{Environment Systems and Decisions} 38 (2018):
  433-457, \url{https://doi.org/10.1007/s10669-018-9708-3}. These
  support possible nonlinear deterioration near saturation, not a fitted
  workplace prediction.} Leah's proposed extension is that such a regime
can leave no action satisfying every live constraint, creating the
conditions in which the fuck-it factor arises.

This yields an uncomfortable but necessary conclusion:

\begin{quote}
\textbf{Ethical quiet quitting may be justified and still be causally
hazardous.}
\end{quote}

Where hidden heroic effort was the final margin, the downstream change
reveals dependence on that effort under those conditions. It need not
identify every upstream cause before the institution acts; those with
greater access and authority bear the duty to investigate them. The
exact consequence may still be difficult to predict. Noncritical backlog
or delay can disclose dependence on withdrawn effort under the observed
conditions; customer, coworker, patient, or public harm cannot be
treated as an honest measurement.

This does not restore an employer's claim to unlimited heroic labor. The
worker should face foreseeable consequences honestly within what they
can reasonably know and control. The institution, acting through those
with the relevant authority, owns the system-level work: priority
ladders, minimum staffing, rate limits, safe-stop conditions, handoffs,
and surge protocols should be specified before the next rush rather than
improvised inside it.

\hypertarget{a-warning-across-unequal-power}{%
\subsection{6.5 A warning across unequal
power}\label{a-warning-across-unequal-power}}

Workers are human beings. Under enough pressure, they may make more
mistakes. Under still more pressure, some will knowingly abandon one
requirement to preserve another. Individual quiet quitting proper should
not disguise retaliation or serious manufactured harm as an accidental
boundary. Collective pressure is a different use of worker power; it
should be named and owned as such. No worker can forecast a complete
cascade from inside a system whose information and controls they do not
possess.

Managers and executives are human beings too, but their power makes the
scale different. The observable managerial analogue appears as ``just
get it done,'' refusal to hear another warning, concealment of bad
numbers, indiscriminate blame, an unsafe labor cut, or an order that
passes an impossible constraint set downward. A private mental state
cannot be diagnosed from outside; the decision and its reach can.
Because managerial decisions scale through other people, one such
decision can become the working conditions of a department, store, or
company.

\begin{quote}
\textbf{Humanness is symmetric. Power is not.}
\end{quote}

A worker's fuck-it event may damage one task or one shift. A manager's
event can establish the conditions under which many workers are
repeatedly forced into the same event. Corporate's version - ``fuck it,
hit the labor number'' - can turn an acute failure of judgment into
standing policy.

The warning is therefore exact:

\begin{quote}
\textbf{A system that demands mutually incompatible standards is not
maintaining all of those standards. It is outsourcing the choice of
which standard will break to the most overloaded person at the worst
moment.}
\end{quote}

The ethical objective is not to demand superhuman adherence after the
fully compliant action-set has already vanished. It is to prevent that
set from vanishing: adequate staffing, explicit priority order, usable
escalation, protected recovery, sufficient equipment, and authority to
slow or stop when the safety floor cannot be preserved.

\hypertarget{employer-size-does-not-settle-the-ethics}{%
\section{7. Employer Size Does Not Settle the
Ethics}\label{employer-size-does-not-settle-the-ethics}}

A large conglomerate may be able to absorb a financial loss, yet
redirect the immediate cost onto a local team, a good supervisor, or
vulnerable coworkers. A small business may be genuinely fragile, yet use
that fragility to claim an unlimited moral entitlement to unpaid labor.

Size matters, but it does not decide the case.

The morally relevant question is where the cascade lands.

Ask:

\begin{enumerate}
\def\labelenumi{\arabic{enumi}.}
\tightlist
\item
  Does the withdrawal expose an existing resource constraint, or create
  a new injury?
\item
  Will the institution absorb the consequence, or merely transfer it
  sideways onto coworkers?
\item
  Has foreseeable or intended third-party pressure been named and
  ethically weighed rather than disguised as accidental?
\item
  Is a decent local manager being made the sole bearer of a failure
  created above their authority?
\item
  Is the withdrawal proportional to the hidden subsidy being removed?
\item
  Does the action improve causal legibility, or merely increase pain?
\item
  Will the withdrawal remove enough slack to push another worker into
  acute overload or a fuck-it event, and can that risk be reduced
  without restoring the hidden subsidy?
\end{enumerate}

For individual quiet quitting proper, the ethical aim is to move
information upward without covertly manufacturing serious suffering
sideways. Collective action can intentionally exert pressure; it should
be named and judged as collective action.

A small business does not own a worker's life because it is fragile. A
conglomerate does not become harmless because it is rich. A worker
should consider foreseeable effects within what they can reasonably know
and control. The institution, with greater information and authority,
bears the larger duty to trace and prevent systemic cascades.

\hypertarget{quiet-quitting-proper-the-last-boundary}{%
\section{8. Quiet Quitting Proper: The Last
Boundary}\label{quiet-quitting-proper-the-last-boundary}}

The following describe the clearest case for quiet quitting proper. They
are a worker's self-audit, not an employer's permission form and not a
fiction that both parties enter the judgment with equal power:

\begin{enumerate}
\def\labelenumi{\arabic{enumi}.}
\tightlist
\item
  \textbf{The problem is persistent and material.} It is not merely one
  disappointing shift or one denied preference.
\item
  \textbf{Structural or managerial expectations exceed the defensible
  job baseline.} The gap may appear as a vague or contradictory written
  role, a material conflict between the role on paper and the job in
  practice, under-resourced volume or pace, incompatible priorities, or
  nominally optional excess enforced through pressure or penalty.
\item
  \textbf{Good-faith repair has been attempted where reasonably safe and
  useful.} The worker has communicated, proposed, clarified, or
  escalated through a viable pathway. Danger, credible retaliation,
  futility, or the absence of a functioning channel can excuse further
  attempts.
\item
  \textbf{The pathways are closed, inert, unreasonably costly, or offer
  only assimilation.} Further attempts predictably produce no repair,
  retaliation, unilateral relabeling of the excess as the job, or more
  uncompensated labor without meaningful access to change.
\item
  \textbf{The excess effort is discretionary and structurally masking.}
  It exceeds the sustainable core role and allows under-resourcing to
  appear adequate.
\item
  \textbf{Core duties remain intact.} Safety, honesty, feasible quality,
  and professional conduct are preserved.
\item
  \textbf{The withdrawal is proportional.} The worker removes the hidden
  subsidy rather than creating an artificial deficit.
\item
  \textbf{For individual quiet quitting proper, avoidable serious harm
  is not covertly manufactured or shifted onto people unable to repair
  the cause.} Collective pressure belongs to a different category.
\item
  \textbf{The purpose is boundary and truth, not punishment.} Anger may
  be present, but revenge does not govern the design.
\item
  \textbf{Foreseeable acute failure modes have been considered within
  the worker's knowledge and control.} Where reasonably safe and useful,
  the worker has warned which duties may become infeasible and requested
  a priority order. Management remains responsible for system-level
  foresight, staffing, and safe-stop design.
\end{enumerate}

At that point, the worker may ethically:

\begin{itemize}
\tightlist
\item
  perform the defined role competently and sustainably;
\item
  decline chronic role creep without compensation, authority, or
  agreement;
\item
  stop volunteering for predictable emergencies produced by normal
  planning;
\item
  maintain agreed availability rather than permanent accessibility;
\item
  refuse impossible combinations and request a priority decision;
\item
  stop taking personal ownership of outcomes determined by managerial
  resource choices;
\item
  conserve time and energy for health, relationships, education, art,
  organizing, job search, or another future;
\item
  remain courteous without manufacturing enthusiasm.
\end{itemize}

Quiet quitting proper is not the abandonment of standards. It is the
narrowing of scope to the standards the institution actually resources.

\hypertarget{the-non-negotiable-invariants}{%
\section{9. The Non-Negotiable
Invariants}\label{the-non-negotiable-invariants}}

For individual quiet quitting proper, the following remain fixed. They
distinguish a personal boundary from concealed retaliation; they are not
a universal tactics code for collective action.

\hypertarget{safety-remains-fixed}{%
\subsection{Safety remains fixed}\label{safety-remains-fixed}}

Do not compromise food safety, physical safety, sanitation, medical
safety, public safety, or any other duty whose failure can seriously
harm another person.

\hypertarget{honesty-remains-fixed}{%
\subsection{Honesty remains fixed}\label{honesty-remains-fixed}}

Do not falsify records, misstate capacity, invent incidents, conceal
urgent defects, or manipulate numbers. A truthful signal is the entire
point.

\hypertarget{core-professional-integrity-remains-fixed}{%
\subsection{Core professional integrity remains
fixed}\label{core-professional-integrity-remains-fixed}}

Do not intentionally produce bad work where competent work remains
feasible within the real role and available resources. Scope may
contract; integrity does not.

\hypertarget{pressure-must-not-be-disguised}{%
\subsection{Pressure must not be
disguised}\label{pressure-must-not-be-disguised}}

Individual quiet quitting should not disguise intentional pressure or
serious manufactured harm as an accidental byproduct. That does not mean
the less powerful must remain harmless enough to be ignored.

\hypertarget{forced-triage-remains-visible}{%
\subsection{Forced triage remains
visible}\label{forced-triage-remains-visible}}

When the full constraint set cannot be satisfied, do not falsely
reclassify the sacrificed requirement as completed, optional, or normal.
Where it is safe and within the worker's knowledge, name what gave way.
Determining the deeper cause and what resources would preserve it next
time belongs primarily to those who control the system. A workaround
that remains invisible can become tomorrow's official capacity
assumption.

\hypertarget{proportionality-remains-fixed}{%
\subsection{Proportionality remains
fixed}\label{proportionality-remains-fixed}}

Withdraw what was being improperly extracted. Do not add retaliation on
top of the withdrawal.

\hypertarget{accountability-remains-fixed}{%
\subsection{Accountability remains
fixed}\label{accountability-remains-fixed}}

A broken institution does not erase responsibility for one's own
actions. Explanation is not exculpation. A boundary must still be owned.

\hypertarget{collective-action-is-a-different-use-of-power}{%
\subsection{Collective action is a different use of
power}\label{collective-action-is-a-different-use-of-power}}

\begin{quote}
\textbf{Do what you feel is right---but do not mistake the feeling for
proof. Test what is true. Know the consequences as far as you can. Own
what you choose.}
\end{quote}

A strike, boycott, slowdown, or coordinated work-to-rule can
intentionally create pressure. Test its factual premises and search for
disconfirmation. Distinguish what is known from what is predicted,
feared, or hoped. Map foreseeable consequences and honest uncertainty;
tell the truth about who chose the action. Owning deliberate pressure is
not self-condemnation, and it does not transfer responsibility for a
brittle or exploitative system onto the worker. Responsibility remains
distributed according to choice, knowledge, power, and control.

\hypertarget{the-ethical-quiet-quitting-test}{%
\section{10. The Ethical Quiet Quitting
Test}\label{the-ethical-quiet-quitting-test}}

Use the questions that bear on the actual situation; no single answer is
a gate. Before quiet quitting proper, ask:

\begin{itemize}
\tightlist
\item
  Have I identified the actual constraint rather than merely named my
  frustration?
\item
  Have I tried the nearest realistic repair pathway where doing so is
  reasonably safe and useful?
\item
  Have I communicated the problem in terms the receiver can act upon?
\item
  Have I proposed at least one feasible intervention or priority choice?
\item
  Have I preserved an honest causal record where safe, lawful, and
  useful?
\item
  Is ``get another job'' a genuinely live option, or only a formal
  possibility that ignores search costs, wage gaps, schedule, benefits,
  transportation, and survival risk?
\item
  Does this workplace offer a credible tangible return for exceptional
  effort - pay, autonomy, mastery, ownership, advancement, belonging, or
  a future I can recognize as mine - or only the identity reward of
  proving I can endure?
\item
  Is the language of the ``hard-working American'' being used to demand
  dream-level sacrifice while the actual job offers no credible pathway
  toward American-Dream goods such as security, ownership, mobility,
  independence, or a self-directed future?
\item
  Am I judging myself against someone whose personal resonance, team,
  timing, or opportunity makes this work a materially different object
  for them?
\item
  Is the effort I plan to withdraw truly discretionary, uncompensated,
  unsustainable, or outside the reasonably defined role?
\item
  Will the withdrawal reveal an existing limit rather than manufacture a
  new failure?
\item
  Am I quiet quitting individually, or choosing a collective pressure
  tactic that should be named and evaluated on its own terms?
\item
  What acute overloads may appear when the hidden buffer is removed?
\item
  Which requirement is most likely to be sacrificed during the next
  rush, and is that requirement ethically eligible for sacrifice?
\item
  Have I requested an explicit priority order and identified a safe-stop
  condition before the fully compliant action-set disappears?
\item
  Are safety, honesty, feasible quality, and coworker welfare protected?
\item
  If this is individual quiet quitting, is its scale and timing
  proportional to the hidden subsidy rather than disguising a different
  tactic?
\item
  For individual quiet quitting proper, am I trying to stop an
  extraction, or trying to make someone suffer?
\item
  Would I be willing to describe the action plainly to a fair-minded
  observer?
\item
  If management repaired the pathway in good faith, would I be willing
  to reconsider?
\end{itemize}

\hypertarget{the-workers-pledge}{%
\section{11. The Worker's Pledge}\label{the-workers-pledge}}

This pledge is not a strict-adherence protocol, checklist, or
employer-facing test. It is a set of reminders: do not self-sacrifice to
preserve institutional appearances; exercise the power you have---or
should have---consciously and ethically; and remain answerable to
reality. Capital-T Truth is the aim, while human understanding is
partial and revisable. Test what you think you know, seek what could
prove you wrong, distinguish knowledge from inference or hope, and
revise when better evidence arrives. Fidelity to reality is not
obedience to authority.

Where a safe and functioning pathway exists, we will try to repair
before we retreat.

We will not lie; when truth-telling is punished, silence is not
falsehood.

We will distinguish a hard day from a structurally impossible demand.

We will not confuse endurance with loyalty, or exhaustion with
excellence.

We will not treat hardship as proof of virtue or another person's
resonance as a universal moral baseline.

We will work hard where effort builds something tangibly worth building,
not merely to prove that we are hard.

We will not confuse the American work ethic with an obligation to keep
feeding a bargain after work has been severed from the dream it was
supposed to make reachable.

We will work to build an owned future, not merely perform the identity
of people who deserve one.

We will not make impossible workloads appear ordinary by consuming
ourselves in private.

We will not mistake assimilation for correction.

We will make those with authority own the choice among incompatible
priorities.

Where we can do so safely, we will name impossible constraint sets
before the fuck-it factor chooses in silence.

When something must give, we will protect the safety floor and make the
sacrificed requirement visible.

We will not treat a shift that merely survived as proof that it was
safely or adequately staffed.

We will not mistake a theoretically available exit for a live one.

We will preserve the connection between resource decisions and their
consequences.

In individual quiet quitting proper, we will not manufacture failure and
misname it as withdrawal.

We will not falsify.

We will not conceal serious safety risks or manufacture serious harm and
call it quiet quitting.

When we act collectively, we will test what is true, name intended
pressure, and own what we choose within what we can reasonably know and
control.

We will give honest labor for honest compensation.

We will help during genuine emergencies without allowing emergency
effort to become the invisible permanent baseline.

We will withdraw the hidden subsidy, not our integrity.

And when every reasonably safe ethical pathway to improvement is closed,
inert, unsafe, or absent, we reserve the right to quiet quit proper.

\hypertarget{a-message-to-managers}{%
\section{12. A Message to Managers}\label{a-message-to-managers}}

Your most capable employees may be hiding your worst process failures.

When a worker stops overfunctioning, the dysfunction may not be new. It
may only have become visible. Do not ask first, ``How do we make this
person care again?'' Ask:

\begin{quote}
What had this person been supplying that our official staffing, process,
and compensation model did not contain?
\end{quote}

A hierarchy that carries commands downward but cannot carry unwelcome
reality upward is not functioning as a communication system. It is
functioning as a performance system in which each level edits the truth
to look competent to the level above.

Repair the channel before demanding renewed investment.

When operative demands exceed the defensible baseline---through
contradiction, role expansion, or impossible volume, pace, staffing,
equipment, or priorities---repair the work system or renegotiate the
expansion. Do not rewrite the description, demand assimilation, and call
that clarity. If a competent employee now looks merely average, the
defined work has not failed; the invisible subsidy has stopped.

Humane restraint is not pathological tolerance. Managers may act against
genuine misconduct. Their obligation is to distinguish it from a
boundary, dissent, organizing, or work-to-rule, and to respond
proportionately rather than treating every loss of discretionary
obedience as betrayal. Greater power requires better discrimination, not
thinner skin.

Ethical quiet quitting proper is a warning case, not a claim that
structurally produced withdrawal always arrives with an explicit ethic
attached. It describes one explicitly articulated attempt to preserve
truth, safety, competent baseline work, and fairness while ending an
unowned subsidy. Other workers may arrive at a similar posture through
the same kind of environment, filtered through their own circumstances
and values, without naming or analyzing it this way. Do not mistake
their lack of a manifesto for lack of cause, and do not let a case of
apathy, refusal, sabotage, or collective pressure elsewhere become a
shortcut for classifying competent baseline work here. Because
management controls classification and discipline, its duty to
distinguish and repair is greater.

Do not mistake your most personally resonant employee for the neutral
human baseline. A worker who loves the mission, trusts the team, sees a
credible future, or has unusually favorable life conditions may freely
invest far more than another competent worker can rationally sustain.
That extra investment is real, but it is contingent. It is not evidence
that everyone else lacks character.

If you want discretionary grind, make the return tangible: pay,
autonomy, skill, ownership, security, credible advancement, meaningful
influence, trustworthy leadership, and a mission people can actually
share. Slogans cannot manufacture resonance, and passion cannot be
converted into permanent staffing capacity.

Do not invoke the ``hard-working American'' while removing the worker's
credible route to the American Dream. If you ask for dream-level
loyalty, flexibility, sacrifice, and output, then wages, security,
advancement, ownership, or meaningful influence must make that sacrifice
intelligible as investment. National mythology cannot repair a severed
causal pathway, and patriotism cannot be booked as compensation.

Do not answer a worker's boundary with ``then leave.'' Continued
attendance may mean that rent, food, medicine, transportation, or health
coverage still depend on the paycheck; it does not prove that the
conditions are acceptable. If the only options you leave are
unsustainable overfunctioning or immediate material insecurity, the
predictable middle state is quiet quitting, not a clean departure.

Reward early warning. Permit explicit priority calls. Distinguish a
worker's boundary from insubordination. Do not punish the person who
makes a hidden constraint measurable. When an employee says that A, B,
and C cannot all be completed safely with the available resources,
management's job is to allocate, not to demand that the contradiction
disappear inside the employee.

Do not model overload as a simple reduction in speed. Near a limit,
ordinary error risk can rise. Beyond it, an employee may have to choose
which requirement gives way while the same pressure makes that tradeoff
harder to judge. If you insist that every target remains mandatory after
the fully compliant action-set is empty, you have not preserved the
standards; you have hidden the selection of the broken standard.

Watch for your own fuck-it factor. We cannot see your private state, but
we can judge the decision and its reach. ``Just get it done'' without a
coherent route is an observable refusal to own the tradeoff. Cutting
labor to hit a number, suppressing a warning to look competent upward,
or blaming a worker because the structural explanation is inconvenient
can impose the same operational constraint collapse at a higher order.
The power difference means your decision can become everyone else's
working conditions.

Quiet quitting proper can therefore expose unanticipated or unwanted
consequences. Treat those effects as evidence relevant to real capacity
and dependence on hidden effort---not as proof of a unique cause, and
not automatically as proof of bad character. Ask what margin
disappeared, which constraint became impossible, and why the
organization had made one person's unsustainable compensation the final
safety system.

The fastest way to produce quiet quitting is to make voice causally
irrelevant while continuing to demand emotional ownership. The fastest
way to turn quiet quitting into a dangerous cascade is to leave every
priority nominally mandatory after the resources needed to satisfy them
have already been removed.

\hypertarget{conclusion-a-last-boundary-not-a-first-weapon}{%
\section{Conclusion: A Last Boundary, Not a First
Weapon}\label{conclusion-a-last-boundary-not-a-first-weapon}}

Quiet quitting should be a last ethical boundary, not a first weapon.

We repair first because work can matter, coworkers matter, customers
matter, and institutions can sometimes learn. Testing what can safely be
tested helps the worker distinguish a repairable problem from a closed
world. It does not require doing management's work or surviving every
channel before a boundary becomes legitimate.

But repair is not a lifelong obligation to bleed into a closed circuit.

When words fail, records are ignored, priorities remain impossible, and
the institution expects private sacrifice to bridge a structural gap it
refuses to acknowledge, the worker may ethically stop paying that
subsidy.

Where exit is not yet a live option, quiet quitting can preserve the
income and recover the bandwidth required to make another future
reachable.

This manifesto is not anti-grind. It defends the right to invest deeply
where effort has somewhere real to go, and the right not to turn being
ground down into proof of identity. Quiet quitting may preserve ambition
by reclaiming it from a system that cannot receive it.

Nor is quiet quitting necessarily a rejection of the American Dream. It
may be evidence that the worker still wants an owned future but no
longer believes this employer connects additional sacrifice to it. The
hard-working-American identity cannot indefinitely survive the
disappearance of the Dream-side pathway. When the bargain becomes all
work and no reachable future, withdrawal is not mysterious; it is
diagnostic.

But the withdrawal is not guaranteed to be mechanically clean. In a
brittle system, removing hidden slack can expose conditions associated
with ordinary error risk and forced triage; in Leah's proposed model,
those conditions can also make the fuck-it factor more likely.
Foreseeable consequences should be faced honestly, but responsibility is
not divided fifty-fifty: it follows choice, knowledge, power, and
control. The worker owns the acts they choose. The institution bears the
larger systemic burden where it designed or maintained the constraint
set. That is not a reason to preserve exploitation. It is a reason to
repair early, specify priorities honestly, and never let one person's
heroics become the only barrier between normal operations and acute
collapse.

We do not manufacture failure. We stop concealing it.

We do not abandon duty. We define it.

We do not withdraw care from human beings. We withdraw unlimited
institutional access to ourselves.

\textbf{Repair first. Make the constraints legible. Withdraw the hidden
subsidy. Then, only when all else fails, quiet quit proper.}

\textbf{We withdraw the subsidy, not the integrity.}

\begin{center}\rule{0.5\linewidth}{0.5pt}\end{center}

\emph{Scope note:} This is an ethical framework, not legal advice.
Employment contracts, collective bargaining agreements, workplace
policies, professional duties, and applicable law may create additional
obligations or protections.

\emph{Evidence and conceptual provenance:} This is a case-anchored
normative manifesto, not a population study. The thesis,
ethical-de-buffering framework, and fuck-it-factor coinage are Leah's
synthesis, developed from lived experience and her TLICA work on agency,
semantic interoperability, configurational action, and the separation of
acts, outcomes, cascades, responsibility, and accountability. Adjacent
research supports parts of the surrounding risk ecology; it does not
validate quiet quitting proper, ethical de-buffering, or the fuck-it
factor as scientific constructs.

\emph{AI-development disclosure:} The author used OpenAI's ChatGPT to
test the argument, identify counterarguments and candidate sources, and
draft and revise portions of the prose. The originating thesis,
lived-experience framing, coined concepts, and editorial direction are
the author's. Publication remains contingent on the author's personal
review of the complete manuscript, independent checking of factual
claims and citations, and approval of the final positions and wording.
Responsibility for any released text remains with the author.

\clearpage
\begingroup
\sloppy
\footnotesize
\setlength{\parskip}{0.08em}
\theendnotes
\endgroup

\end{document}
